\section{Tate--Shafarevich groups of cyclotomic twists: the window bound}\label{sec:window}

\providecommand{\windowQZ}{\Q_\ell/\Z_\ell}
\providecommand{\windowcyc}{\chi_{\mathrm{cyc}}}
\providecommand{\windowX}{\mathcal{X}}
\providecommand{\windowH}{\mathcal{H}}
\providecommand{\windowHf}{H^1_{f}}
\providecommand{\windowsh}{\mathrm{shaLen}}

This section is devoted to the proposition $\sfname{VWin}$ of \cref{def:vwin} and to its proof from the
statement $\sfname{Row0b}$ on Galois representations attached to polarized regular algebraic cuspidal
automorphic representations (\cref{def:row0b}), which is a theorem by \cref{thm:row0b}. $\sfname{VWin}$ is an
explicit upper bound for the orders of the groups
\[
  \Sha^1\bigl(G_{E,S_\ell},\windowQZ(-j)\bigr)\ \subset\ H^1\bigl(\OO_E[1/\ell],\windowQZ(-j)\bigr)
\]
of classes unramified outside $\ell$ and locally trivial above $\ell$, simultaneously for all primes $\ell$ and all
twists $1\le j<x^{401}$, where $E$ runs through the abelian CM fields $E_{\psi,m,q,d}$ underlying the host fields
of \cref{sec:hosts}, with $2^m\asymp x^{2+8\eps}$, and $x$ lies in a suitable window $[Y,Y^2]$. In \cref{sec:top}
the bound supplies host fields at all scales of such windows (\cref{thm:supply}); the logarithm of the product of
the orders is the quantity entering the Sha link among the links of \cref{def:links}.

\emph{Plan.} $\Sha$ is split into its $2$-part and, for odd $\ell$, the eigenspaces $\Sha^\pm$ of complex
conjugation (\cref{sec:window:statement,sec:window:arch}). The minus part is bounded by generalized Bernoulli
numbers, through a bound of main-conjecture strength proved with an Euler system of Gauss sums and the
reciprocity law (\cref{sec:window:minus}). The $2$-part and the finitely many bad primes are treated by Iwasawa
theory in cyclotomic $\Z_\ell$-towers and an explicit bound for $H^2$; the window is chosen by a pigeonhole
argument (\cref{sec:window:towers}). The plus part is reduced to Dirichlet characters over $\Q$ and bounded, case
by case, by Eisenstein congruence numbers for Galois representations of $\Uni(3)$- and $\GSp_4$-type
(\cref{sec:window:plus,sec:window:eis}), using a vanishing theorem for Bloch--Kato Selmer groups proved by an
Euler system argument (\cref{sec:window:kato}). All inputs, including the base change and descent theorems used,
are proved in the formal development \cite{LeanSource2026} (see also \cref{sec:literature}). The reduction of
$\sfname{VWin}$ to them, which occupies roughly half a million lines there, is new; we state its intermediate
results precisely and outline their proofs. The minus part follows the classical Euler-system method closely; the
bounds for the plus parts by Eisenstein congruence numbers are the main new ingredient.

% ---------------------------------------------------------------------------
\subsection{The groups and the statement}\label{sec:window:statement}

\begin{definition}[Tate--Shafarevich groups of Tate twists]\label{def:window:sha}
Let $F$ be a number field, $\Omega$ an algebraic closure, $G_F=\Gal(\Omega/F)$ and $\ell$ a prime. Let $S_\ell$ be
the set of primes of $F$ above $\ell$ and $G_{F,S_\ell}$ the quotient of $G_F$ by the closed normal subgroup
generated by the inertia groups of the places of $\Omega$ above finite primes $v\notin S_\ell$, i.e.\ the Galois
group of the maximal extension unramified at all finite primes outside $\ell$ (no condition at infinity). An
\emph{$\ell$-datum} $D$ is a choice of a place of $\Omega$ above each $v\in S_\ell$, with decomposition group
$D_v\subset G_F$. For $j\in\Z$ let $\windowQZ(-j)$ be the discrete module on which $G_{F,S_\ell}$ acts through
$\windowcyc^{-j}$, and let $\Sha_D(j)$ be the kernel of the restriction map
$H^1(G_{F,S_\ell},\windowQZ(-j))\to\prod_{v\in S_\ell}H^1(D_v,\windowQZ(-j))$ (continuous cochains). Put
\[
  \windowsh(D,j):=\len_{\Zl}\Sha_D(j)\in\N\cup\{\infty\},\qquad
  \mathrm{VLen}(D,j):=\len_{\Zl}H^2\bigl(G_{F,S_\ell},\Zl(j+1)\bigr);
\]
thus $\#\Sha_D(j)=\ell^{\windowsh(D,j)}$ when this is finite.
\end{definition}

For $\ell$ odd or $F$ totally imaginary, $\Sha_D(j)$ is the usual $\Sha^1(G_{F,S_\ell},\windowQZ(-j))$ in the
\'etale cohomology of $\OO_F[1/\ell]$; all fields below are CM, so the omission of the real places is
immaterial. Localizing at all places above $\ell$ gives the same group, so $\windowsh(D,j)$ does not depend on
$D$; we write $\windowsh(F,\ell,j)$ when the choice is irrelevant.

\begin{remark}[What $\windowsh$ measures]\label{rem:window:meaning}
The group $\Sha^1(G_{F,S_\ell},\windowQZ(-j))$ is, by Poitou--Tate duality \cite{NeukirchSchmidtWingberg2008,Milne2006},
Pontryagin dual to $\Sha^2(G_{F,S_\ell},\Zl(1+j))\subset H^2(\OO_F[1/\ell],\Zl(1+j))$, which is finite for $j\ge1$
\cite{Soule1979}. For $j\ge1$ this is the $\ell$-part of the \'etale wild kernel $WK^{\mathrm{et}}_{2j}(F)$
\cite{Banaszak1993} (for $j=1$, of the wild kernel in $K_2(F)$ \cite{Tate1976}), hence, by the Quillen--Lichtenbaum
conjecture \cite{Voevodsky2011}, the $\ell$-part of a wild kernel in $K_{2j}(\OO_F)$; so
$\sum_\ell\log\ell\cdot\windowsh(F,\ell,j)=\log\#WK^{\mathrm{et}}_{2j}(F)$. For abelian $F$ these orders are given by
the Tamagawa number conjecture for Dirichlet motives \cite{BlochKato1990}, proved for odd $\ell$ by Huber--Kings
\cite{HuberKings2003}: on the components where $L(\theta,-j)\ne0$ (the minus components below) by generalized
Bernoulli numbers, compare \cite{MazurWiles1984,CoatesSinnott1974}; on the others (plus components) by leading terms
involving Borel regulators. The groups bounded here are thus wild kernels in $K$-theory, not Bloch--Kato Selmer
groups of motives of non-zero weight.
\end{remark}

\begin{definition}[Plus and minus parts]\label{def:window:pm}
Let $F$ be CM with complex conjugation $c$, and $j\ge0$. Then $c$ acts on $H^1(G_{F,S_\ell},\windowQZ(-j))$ and
preserves $\Sha_D(j)$. Let $\Sha^\pm_D(j)=\Sha_D(j)\cap\ker\bigl(c\mp(-1)^{j}\bigr)$, with $\Zl$-lengths
$\mathrm{plusLen}(D,j)$ and $\mathrm{minusLen}(D,j)$.
\end{definition}

For $F$ abelian over $\Q$, the $\theta$-component for a character $\theta$ of $\Gal(F/\Q)$ lies in $\Sha^-$ if and
only if $\theta(-1)=(-1)^{j+1}$, the parity for which $L(\theta,-j)\ne0$.

\begin{lemma}\label{lem:window:pm}
If $\ell\ne2$, then $\windowsh(D,j)=\mathrm{plusLen}(D,j)+\mathrm{minusLen}(D,j)$.
\end{lemma}

\begin{proof}
As $2\in\Zl^\times$ and $c^2=1$, the idempotents $\frac12(1\pm(-1)^jc)$ split $\Sha_D(j)$; length is additive.
\end{proof}

\begin{definition}[The fields $E_{M,m,q,d}$]\label{def:window:fields}
An \emph{abelian datum} $M=(m_M,H_M)$ is a modulus $m_M$ with a subgroup $H_M\subset(\Z/m_M)^\times$ of index
$\deg M$. For $m\ge0$ and primes $q,d$, a number field $E$ is an \emph{$E_{M,m,q,d}$} if for some $N$ divisible
by $2^{m+2}$, $m_M$, $q$ and $d$ there is an embedding $E\hookrightarrow\Q(\zeta_N)$ whose image is the fixed field of
\[
 \Bigl\{a\in(\Z/N)^\times:\ a\equiv\pm1\ (\mathrm{mod}\ 2^{m+2}),\ \ a\bmod m_M\in H_M,\ \
 \Bigl(\frac aq\Bigr)=\Bigl(\frac ad\Bigr)=1\Bigr\}.
\]
For a primitive character $\psi$ of conductor $N_\psi$ and order $o_\psi$ we write $E_{\psi,m,q,d}$ for
$E_{M_\psi,m,q,d}$, where $M_\psi=(N_\psi,\ker\psi)$ and $\deg M_\psi=o_\psi$ (formally $\sfname{IsChiField}$).
Thus $E_{\psi,m,q,d}$ is the compositum of $\Q(\zeta_{2^{m+2}})^+$, the field cut out by $\psi$ and the quadratic
fields cut out by the two Jacobi symbols.
\end{definition}

\begin{lemma}\label{lem:window:cm}
If $q\equiv3\pmod4$, every $E_{\psi,m,q,d}$ is a CM field, and $[E:\Q]\le\cond(E)\le2^{m+2}N_\psi qd$.
\end{lemma}

\begin{proof}
$E$ is abelian and contains the Gauss sum $G$ with $G^2=-q$, which is fixed by the defining subgroup; so
$\Q(\sqrt{-q})\subset E$, and an abelian totally imaginary field is CM. The embedding descends to $\Q(\zeta_{N'})$
with $N'=\lcm(2^{m+2},N_\psi,q,d)$.
\end{proof}

We recall $\sfname{VWin}$ (\cref{def:vwin}), writing out all quantifiers.

\begin{definition}[The window bound]\label{def:window:vwin}
Let $\eps>0$.
\begin{enumerate}[label=(\alph*)]
\item $\sfname{PageN}(\eps,x,m)$: $x^{2+8\eps}\le2^m\le2x^{2+8\eps}$.
\item $\sfname{WinPair}(\eps,Y;q,d)$: $q\equiv3\pmod8$ and $d\equiv7\pmod{16}$ are primes,
  $\frac12Y^{\eps/5}<\log q\le Y^{\eps/5}$, $\frac12Y^{\eps/5}(\log Y)^2<\log d\le Y^{\eps/5}(\log Y)^2$,
  $q\mid d+1$, and every odd prime $p\le Y^{\eps/5}$ divides $d+1$.
\item $\sfname{VVle}(E,t,B)$: for every algebraic closure of $E$ and every choice of $p$-data $D_p$ for all
  primes $p$, the lengths $\windowsh(D_p,j)$, $0<j<t$, are finite and, for every finite set $S$ of primes,
  \[\textstyle\sum_{\ell\in S}\log\ell\sum_{j=1}^{t-1}\windowsh(D_\ell,j)\le B,\quad\text{that is,}\quad
  \log\prod_{\ell}\prod_{j=1}^{t-1}\#\Sha_{D_\ell}(j)\le B.\]
\item $\sfname{VWinAt}(\psi,C_t)$: for every $\eps\in(0,\frac1{52})$ and every $m:\R\to\N$ with
  $\sfname{PageN}(\eps,x,m(x))$ for all $x\ge1$, there are $J\ge1$ and $X_0$ such that for all $X\ge X_0$ and all
  $(q,d)$ with $\sfname{WinPair}(\eps,X^{1/J};q,d)$ there is $Y\in[X^{1/J},X^{1/2}]$ with the following property:
  for all $x\in[Y,Y^2]$, every $E_{\psi,m(x),q,d}$ satisfies
  $\sfname{VVle}\bigl(E,\lfloor x^{C_t}\rfloor,\sqrt{\exp(x^{\eps})}\bigr)$.
\item $\sfname{VWin}$: $\sfname{VWinAt}(\psi,401)$ holds for every primitive character $\psi$.
\end{enumerate}
\end{definition}

\begin{theorem}[Window bound]\Fmark\label{thm:vwin}
$\sfname{Row0b}$, with the local parameter predicate $\sfname{IsRecT}$ of \cref{sec:galois}, implies
$\sfname{VWin}$. Consequently, by \cref{thm:row0b}, $\sfname{VWin}$ holds.
\formal{UpperTheory.One.vwin_of_row}
\end{theorem}
%P2PTR thm:vwin
\begin{proof}See \cref{sec:pf:vwin}.
\end{proof}


\begin{remark}[On the statement]\label{rem:window:statement}
For $x$ in a window, the degree $[E:\Q]$ has order $x^{2+8\eps}$, one has $\log\cond(E)\ll x^{\eps/5}(\log x)^2$,
and the twists satisfy $j<x^{401}$. By the conjectures of Lichtenbaum and Bloch--Kato one expects the logarithm of
$\#WK^{\mathrm{et}}_{2j}(E)$ to be polynomial in $j$,
$[E:\Q]$ and $\log|d_E|$, so the bound $\exp(x^\eps/2)$ is generous; the content lies in the finiteness for all $\ell$
and $j$, the uniformity, and above all the bound for the plus parts, whose analytic counterpart would involve lower
bounds for Borel regulators; such bounds do not seem to have been available before. The bound is claimed only for
$x$ in some window $[Y,Y^2]$ with $X^{1/J}\le Y\le X^{1/2}$, produced by \cref{lem:window:window}. The existence of
pairs $(q,d)$ with $\sfname{WinPair}$ is not part of $\sfname{VWin}$; such pairs are produced in \cref{sec:top},
where the exponent $401$ is used as $400$.
\end{remark}

% ---------------------------------------------------------------------------
\subsection{Architecture of the proof}\label{sec:window:arch}

The formal proof first assembles $\sfname{Row0b}$, a formal proof of base change from definite unitary groups and
formally proved results from the literature into a package $\sfname{LiteratureRankNHypothesis}$: representation
theory of $p$-adic $\GL_n$; automorphic forms on $\GL_n$ \cite{JacquetShalika1976,JacquetShalika1981,Clozel1990};
Galois cohomology (Poitou--Tate and local Tate duality, Euler characteristics, orthogonality of unramified classes,
the Greenberg--Wiles formula \cite{Wiles1995}, Bloch--Kato's local theory \cite{BlochKato1990}); and complex
multiplication \cite{Shimura1971,DeligneRibet1980}. A larger package $\sfname{HypothesesH'}$ adds further results of
Clozel \cite{Clozel1990} and Borel \cite{Borel1983}, Galois representations for Hilbert modular forms
\cite{Carayol1986,Taylor1989}, strong multiplicity one, definite quaternion algebras \cite{Vigneras1980}, and the
statements of the following remark.

\begin{remark}[Base change and descent]\label{rem:window:bc}
The following known theorems are used. Each is stated as a proposition and given a formal proof in the
development, by arguments which differ from the published ones; only the statements matter here.
\begin{enumerate}[label=(\roman*)]
\item Weak base change from definite unitary groups attached to CM extensions $E/E^+$ unramified at all finite
  places to isobaric representations of $\GL_n(\A_E)$, compatible at split places, at unramified inert places and
  at infinity \cite{Labesse2011}.
\item Descent of conjugate self-dual regular algebraic cuspidal representations of $\GL_n(\A_E)$, $2\le n\le4$, to
  definite unitary groups with prescribed split components, a statement of the type proved in \cite{Labesse2011}.
\item Cyclic base change and descent of prime degree for $\GL_n$, and automorphic induction for quadratic
  extensions \cite{ArthurClozel1989}.
\item Quadratic base change for $\GL_2$, and descent to $\Q$ of cuspidal regular algebraic representations of
  $\GL_2$ over imaginary quadratic fields whose Satake parameters are Galois invariant at almost all places
  \cite{Langlands1980}.
\end{enumerate}
\end{remark}

\begin{theorem}[Main reduction]\Fmark\label{thm:window:HE}
$\sfname{HypothesesH'}$ implies $\sfname{VWin}$.
\formal{UpperTheory.V.HV.DescentVanishingReductions.vwin_of_hypotheses_H_E_alt}
\end{theorem}
%P2PTR thm:window:HE
\begin{proof}See \cref{sec:pf:vwin-he}.
\end{proof}


The proof combines the following theorems of the development: $\sfname{LemmaK}$ (\cref{thm:window:lemmaK}), the
minus part (\cref{prop:psibound}), the $2$-part and the bad layers (\cref{thm:window:bab}), and the plus part
(\cref{thm:plus-part}), whose inputs $\sfname{CountI}$, $\sfname{CountII}$, $\sfname{CountIII}$, $\sfname{CountV}$
come from automorphy statements of $R=T$ type, $\sfname{LemmaK}$ and \cref{thm:kato} (\cref{sec:window:eis}), and
$\sfname{CountZ}$ from \cref{prop:window:propZ}; the last step is \cref{prop:window:assembly}. \Cref{thm:vwin}
follows by short implications: $\sfname{Row0b}$, (i) of \cref{rem:window:bc} and the literature results give
$\sfname{LiteratureRankNHypothesis}$; with (ii) they give $\sfname{HypothesesH'}$; and \cref{thm:window:HE} applies.

\begin{definition}[Partial bounds]\label{def:window:partial}
$\sfname{VVat}(\ell,E,t,B)$: for every $\ell$-datum $D$ of $E$ the lengths $\windowsh(D,j)$, $0<j<t$, are finite
with sum $\le B$. For $E$ CM, $\sfname{VVminusLe}(E,t,B)$ and $\sfname{VVplusLe}(E,t,B)$ are defined as
$\sfname{VVle}(E,t,B)$ with $\mathrm{minusLen}$, resp.\ $\mathrm{plusLen}$, in place of $\windowsh$ and with $\ell$
restricted to odd primes.
\end{definition}

By \cref{lem:window:pm}, if $E$ satisfies $\sfname{VVat}(2,E,t,B_2)$, $\sfname{VVminusLe}(E,t,B_-)$ and
$\sfname{VVplusLe}(E,t,B_+)$, then it satisfies $\sfname{VVle}(E,t,B)$ with $B=B_2\log2+B_-+B_+$. Put
\[
 \beta(C,o,n,N,t)=3nt^2\log(tN)+C\sum_{\ell\mid o,\ \ell\ne2}\ell\,n^2t^2\log(t\ell N)^2 .
\]

\begin{proposition}[Final assembly]\Fmark\label{prop:window:assembly}
Let $C_t\ge1$, and assume:
\begin{enumerate}[label=(\alph*)]
\item ($\sfname{ThmBcFree}$) there is $C_-$ such that for every abelian datum $M$, all $k,t\ge2$ and all primes
  $q\equiv3\bmod8$, $d\equiv7\bmod16$ prime to $m_M$, every CM field $E=E_{M,k,q,d}$ satisfies
  $\sfname{VVminusLe}(E,t,B)$ with $B=\beta(C_-,\deg M,[E:\Q],\cond E,t)$;
\item (joint window) for every primitive $\psi$, $\eps\in(0,\frac1{52})$, $A_1>0$ and all $m,t:\R\to\N$ with
  $2\le m(x)$, $2^{m(x)}\le x^{A_1}$, $2\le t(x)\le x^{C_t}$, there are $J,C,P$ such that for all large $X$ and all
  primes $q\equiv3\bmod8$, $d\equiv7\bmod16$ prime to $N_\psi$ with $q,d>\max(2X^{C_t}+1,P)$ and
  $\log(qd)\le X^{\eps/(4J)}$, there is $Y\in[X^{1/J},X^{1/2}]$ such that for $x\in[Y,Y^2]$ every
  $E=E_{\psi,m(x),q,d}$ satisfies $\sfname{VVat}(2,E,t(x),e^{x^{4\eps/5}})$ and $\sfname{VVplusLe}(E,t(x),B_+)$,
  where $B_+=(\cond(E)t(x))^C+Ce^{x^{4\eps/5}}$.
\end{enumerate}
Then $\sfname{VWinAt}(\psi,C_t)$ holds for every primitive $\psi$.
\formal{UpperTheory.V.vwinAt_of_joint}
\end{proposition}

\begin{proofoutline}
Take $t(x)=\lfloor x^{C_t}\rfloor$. Elementary estimates show: window pairs at scale $X^{1/J}$ satisfy the
conditions in (b) for large $X$; $\log(\cond(E)t)\le2x^{\eps/4}$, by \cref{lem:window:cm} and $\sfname{PageN}$;
$\beta(C_-,o,n,N,t)\le K'e^{8W}$ if $n\le N$ and $\log(Nt)\le2W$; and
$e^{x^{4\eps/5}}\log2+K'e^{8x^{\eps/4}}+u^C+Ce^{x^{4\eps/5}}\le\sqrt{\exp(x^{\eps})}$ for large $x$ and
$\log u\le2x^{\eps/4}$. Adding the three partial bounds gives $\sfname{VVle}$.

%P2PTR prop:window:assembly
A complete proof is given in \cref{sec:pf:vwin}.
\end{proofoutline}

Hypothesis (a) is \cref{prop:psibound}(ii) (in the formal statement it appears in an equivalent form adapted to
the pairs $(m,t)$); hypothesis (b) follows from \cref{thm:window:bab,thm:plus-part}.

% ---------------------------------------------------------------------------
\subsection{The minus part}\label{sec:window:minus}

For a Dirichlet character $\chi$ with values in $\Qlbar$ let $B_{k,\chi}$ be the generalized Bernoulli numbers of
the primitive character attached to $\chi$; for $\chi$ modulo $N$ put
$\alpha_k(\chi)=N\cdot\mathrm{primorial}(k+1)\cdot B_{k,\chi}$, where $\mathrm{primorial}(n)$ is the product of the
primes $\le n$. Recall that $L(\chi,1-k)=-B_{k,\chi}/k$.

\begin{definition}[$\sfname{PsiBound}$]\label{def:window:psibound}
$\sfname{PsiBound}$ is the statement: let $E\subset\Q(\zeta_N)$ be an abelian CM field, characters of $\Gal(E/\Q)$
being viewed as Dirichlet characters modulo $N$; let $\ell\nmid2[E:\Q]$ be a prime, $\OO$ the ring of integers of a
finite extension of $\Ql$, $j\ge1$, and $\theta:\Gal(E/\Q)\to\OO^\times$ a character with $\theta(-1)=(-1)^{j+1}$.
Then
\[
 \len_\OO\Bigl(\Sha^1\bigl(G_{E,S_\ell},\OO\otimes\windowQZ(-j)\bigr)[\theta]\Bigr)\ \le\ v_\OO\bigl(\alpha_{j+1}(\theta)\bigr),
\]
where $[\theta]$ denotes the $\theta$-eigenspace for the action of $\Gal(E/\Q)$.
\end{definition}

The parity condition gives $L(\theta,-j)\ne0$, so the bound is finite. It has the strength predicted by the main
conjecture \cite{MazurWiles1984,HuberKings2003} for each minus component, with the explicit extra factor
$N\cdot\mathrm{primorial}(j+2)$.

\begin{proposition}[Minus part]\Fmark\label{prop:psibound}
\begin{enumerate}[label=(\roman*)]
\item Poitou--Tate duality and the orthogonality of unramified classes under local duality imply
  $\sfname{PsiBound}$.\formal{UpperTheory.V.ML.iv}
\item $\sfname{PsiBound}$ implies hypothesis (a) of \cref{prop:window:assembly}.\formal{UpperTheory.V.Odd.thmBcFree_of}
\end{enumerate}
\end{proposition}
%P2PTR prop:psibound
\begin{proof}See \cref{sec:pf:vwin}.
\end{proof}


The proof of (i) is an Euler-system argument with Gauss sums in the style of Thaine, Kolyvagin and Rubin
\cite{Thaine1988,Kolyvagin1990,Rubin1990}, transposed to Tate twists; global duality enters only through the
following lemma. For $\alpha:G_\Q\to\OO^\times$ we write $\alpha^*=\windowcyc\alpha^{-1}$.

\begin{lemma}[Reciprocity at an auxiliary prime]\Fmark\label{lem:window:rec}
Let $\ell$ be odd, $0\ne a\in\OO$, $w_0$ a generator of the $a$-torsion of $\OO\otimes\windowQZ$, and
$\alpha:G_\Q\to\OO^\times$ with $\alpha\equiv1\bmod a$ on an open subgroup. Let $\ell\mid N$, let $q\ne\ell$ be
prime, and let $c$, $n$ be cocycles for $\alpha$, $\alpha^*$ with values killed by $a$ such that: at primes
$p\nmid Nq$, $\alpha\equiv1\bmod a$ on inertia and $c$, $n$ vanish on inertia; at primes $p\mid N$, $p\ne q$, $c$ is
locally a coboundary; at a place over $q$, $\alpha\equiv\alpha^*\equiv1\bmod a$ on the decomposition group and $c$
vanishes on inertia. If $\tau$ lies in the inertia group at $q$ and $n(\tau)=y\,w_0$, then $y\,c(g)=0$ for all $g$
in the decomposition group at $q$.
\end{lemma}

\begin{proof}[Sketch]
The local invariants of $c\cup n$ vanish at primes $p\nmid Nq$ (unramified orthogonality), at $p\mid N$ ($c$ is
locally trivial) and at infinity; by the reciprocity law, a consequence of Poitou--Tate duality, the invariant at
$q$ vanishes too. There both characters are trivial modulo $a$, $c$ is unramified and $n$ is tamely ramified, and
the local pairing is $c(\Frob_q)\cdot y$ up to a unit.

%P2PTR lem:window:rec
A complete proof is given in \cref{sec:pf:vwin}.
\end{proof}

\begin{proof}[Outline of the proof of \cref{prop:psibound}]
(i) \emph{Reductions.} Since $\ell\nmid[E:\Q]$, inflation--restriction identifies the $\theta$-component with
$\Sha_\ell(\Q,\OO\otimes\windowQZ(\alpha))$, $\alpha=\theta^{-1}\windowcyc^{-j}$, the group of classes of $G_\Q$
vanishing on all inertia groups and on the decomposition groups above $\ell$. One may pass to the primitive
character, since $\alpha_{j+1}$ at the conductor divides $\alpha_{j+1}$ at any level. If $\alpha\equiv\windowcyc$
modulo the maximal ideal on all of $G_\Q$, the group vanishes by a direct computation. Otherwise fix
$F=f'\ell^m$ ($f'$ the prime-to-$\ell$ conductor, $m$ large) such that $\alpha\bmod\ell^m$ factors through
$(\Z/F)^\times$ and is odd.

\emph{Stickelberger values.} For $b$ prime to $F$ let $\vartheta_b(t)=\lfloor bt/F\rfloor-\lfloor b(F-t)/F\rfloor$
and $z_b=\sum_{t\in(\Z/F)^\times}\vartheta_b(t)\alpha(\sigma_t)^{-1}$ ($\sigma_t$ lifting $t$), the image of a
Stickelberger element. A congruence between such values and Bernoulli numbers shows that for every $a_0\ne0$ one
can choose $m$ and $b$ with $a_0z_b\mid\ell^m$ and $z_b\mid\alpha_{j+1}(\theta)$.

\emph{An Euler system of Gauss sums.} Fix a prime $r\equiv1\pmod F$, and call $q$ admissible if $q\equiv1\pmod F$,
$q\ne r$ and $\ell\nmid\ord(r\bmod q)$. By Stickelberger's theorem, applied through Jacobi sums, there are elements
$X_d\in\Q(\zeta_{Fd})$, for squarefree products $d$ of admissible primes, with $X_d\bar X_d=1$, which are units away
from $r$, satisfy Euler-system norm and congruence relations, and such that $X_1$ generates $\mathfrak r^{\vartheta_b}$
for a prime $\mathfrak r$ above $r$. Kolyvagin derivatives, Kummer theory and projection to the
$\alpha^{-1}$-eigenspace give classes $\nu(L)$
for $\alpha^*$, indexed by finite sets $L$ of admissible primes, killed by $\ell^m$ and unramified outside
$Fr\prod_{q\in L}q$, such that: the inertia values of $\nu(\emptyset)$ at $r$ have annihilator $\ell^m/z_b$; and for
$q\notin L$ the inertia values of $\nu(L\cup\{q\})$ at $q$ are unit multiples of the value of $\nu(L)$ at a
Frobenius element at $q$.

\emph{Kolyvagin's induction.} Let $S(L)$ be the group of cocycles of the Sha-type group killed by $a_0$ which are
locally trivial at $r$ and at the primes of $L$, and $a_L$ the annihilator of the values of $\nu(L)$. If
$S(L)\ne0$, pick $0\ne c\in S(L)$; Chebotarev's theorem (using $\alpha\not\equiv\windowcyc$) gives an admissible $q$
with a Frobenius element $\varphi$ such that $c(\varphi)\ne0$ and $\nu(L)(\varphi)$ generates the values of $\nu(L)$.
By \cref{lem:window:rec} for $c$ and $n=\nu(L\cup\{q\})$, localization at $q$ maps $S(L)$, with kernel
$S(L\cup\{q\})$, to a cyclic group killed by $e$, where $a_{L\cup\{q\}}=e\,a_L$ and $e$ is not a unit. Hence
$\len S(L)+v(a_L)\le\len S(L\cup\{q\})+v(a_{L\cup\{q\}})$, where $v(a_L)$ increases strictly and is bounded by
$v(\ell^m)$. Iterating from $L=\emptyset$ and using the ramification of $\nu(\emptyset)$ at $r$ bounds the length of
the $a_0$-torsion by $v(z_b)\le v(\alpha_{j+1}(\theta))$, for every $a_0$.

(ii) For odd $\ell\nmid[E:\Q]$, summing $\sfname{PsiBound}$ over $\theta$ and estimating products of normalized
Bernoulli numbers gives $\sum_\ell\log\ell\cdot\mathrm{minusLen}(D_\ell,j)\le[E:\Q]\log\bigl(16N^{j+2}(j+1)!\bigr)$. For
odd $\ell\mid[E:\Q]$ (then $\ell\mid\deg M$), annihilation of $\Sha^-$ by higher Stickelberger elements, in the manner
of Coates--Sinnott \cite{CoatesSinnott1974}, and a bound for the number of generators give
$\mathrm{minusLen}(D_\ell,j)\log\ell\le C\ell[E:\Q]^2(j+1)\log(j\ell N)^2$. Summing over $1\le j<t$ gives $\beta$.
\end{proof}

For $j=0$ the argument is Rubin's Euler-system proof of the bound for the minus parts of class groups of cyclotomic
fields by Stickelberger elements \cite{Rubin1990}, which together with the analytic class number formula gives the
main conjecture of Mazur--Wiles. Here only the upper bound is needed, and Iwasawa theory does not enter.

% ---------------------------------------------------------------------------
\subsection{Iwasawa towers, Lemma K and the choice of the window}\label{sec:window:towers}

\begin{definition}[$\sfname{LemmaK}$]\label{def:window:lemmaK}
$\sfname{LemmaK}(C_K)$: for every prime $\ell$, every totally imaginary number field $L$ containing $\mu_\ell$
($\mu_4$ if $\ell=2$), every $\ell$-datum $D$ of $L$ and every $j\ge1$, the group $H^2(G_{L,S_\ell},\Zl(j+1))$ is finite
and $\log\#H^2(G_{L,S_\ell},\Zl(j+1))\le\exp\bigl((2j[L:\Q])^{C_K}\log(2|d_L|)\bigr)$.
\end{definition}

\begin{theorem}[Lemma K]\Fmark\label{thm:window:lemmaK}
$\sfname{LiteratureRankNHypothesis}$ implies $\sfname{LemmaK}(C_K)$ for some $C_K$.
\formal{UpperTheory.V.ZK.lemmaK_litN_alt}
\end{theorem}
%P2PTR thm:window:lemmaK
\begin{proof}See \cref{sec:pf:vwin}.
\end{proof}


The bound is doubly exponential but explicit and uniform in $\ell$. Finiteness is Soul\'e's theorem
\cite{Soule1979}, which rests on Borel's computation of the ranks of $K$-groups of rings of integers
\cite{Borel1974}; correspondingly, the formal proof uses Poitou--Tate and local duality together with a rank bound
of Borel type in the stable range, itself given a formal proof, and makes every step quantitative.

\begin{littheorem}[Ferrero--Washington {\cite{FerreroWashington1979}}]\label{lit:window:fw}
Let $K$ be an abelian number field containing $\mu_\ell$ ($\mu_4$ if $\ell=2$). There is $r$ such that the
$\ell$-rank of the class group of $K(\mu_{\ell^n})$ is at most $r$ for all $n$.
\end{littheorem}

This is the vanishing of Iwasawa's $\mu$-invariant for abelian fields, formally proved in the development. It
implies that for such $L$ the Galois group $X'_\infty$ of the maximal unramified abelian pro-$\ell$ extension of
$L(\mu_{\ell^\infty})$ in which the primes above $\ell$ split completely is a finitely generated $\Zl$-module.

\begin{definition}[Tower bounds]\label{def:window:abound}
Let $L_0$ be a number field, $\ell$ a prime, $R\in\N$ and $j_0=\lfloor\log_\ell R\rfloor+3$. A \emph{layer} $L_1/L_0$
of degree $\ell^k$ is a field $L_0(\mu_{\ell^n})$ of degree $\ell^k$ over $L_0$. $\sfname{ABound}(\ell,R,L_0)$ means:
there are a complete non-archimedean field $K\supset\Ql$, $c\in\N$ with $\mu_{\ell^c}\subset L_0$ and points
$s_1,\dots,s_{R'}\in K$, $R'\le R$, such that for all $j\ge1$, all layers $L_1/L_0$ of degree $\ell^k$ and all
$B\in\N$ with $|s_i-j|\ge\ell^{-B}$ for all $i$,
\[
 \windowsh(L_1,\ell,j)\le R'\bigl(\max(k,j_0)+c+B\bigr);
\]
moreover $\windowsh(L_1,\ell,j)\le\mathrm{VLen}(L_2,j)+R(k-j_0)$, where $L_2$ is the layer of degree $\ell^{j_0}$.
\end{definition}

\begin{lemma}\Fmark\label{lem:window:tower}
Assume Poitou--Tate duality and local finiteness and duality. If $L_0$ is totally imaginary, contains $\mu_\ell$ ($\mu_4$ if $\ell=2$), and
$X'_\infty$ is generated by $R\ge1$ elements over $\Zl$, then $\sfname{ABound}(\ell,R,L_0)$ holds.
\end{lemma}
%P2PTR lem:window:tower
\begin{proof}See \cref{sec:pf:vwin}.
\end{proof}


Over a layer, $\Sha$ is controlled by twisted coinvariants of $X'_\infty$; the points $s_i$ play the role of the
$\ell$-adic zeros, in the variable $j$, of its characteristic power series, and away from them the groups grow at
most linearly along the tower ($\mu=0$, $\lambda\le R$). Twists $j$ which are $\ell$-adically very close to some $s_i$
are \emph{exceptional}; for them only the second bound is available, made explicit by $\sfname{LemmaK}$ on the
fixed layer $L_2$. Exceptional twists are avoided by the choice of the window.

\begin{lemma}[Window lemma]\label{lem:window:window}
Let $R\in\N$, $0<\eta\le C_t$, $X>1$, and let $\mathrm{Exc}$ be a set of at most $R$ positive integers. Put
$J=2^{R+2}(C_t/\eta)^R$. There is $Y\in[X^{1/J},X^{1/2}]$ such that for every $x\in[Y,Y^2]$, every
$e\in\mathrm{Exc}$ with $e\le x^{C_t}$ satisfies $e\le x^{\eta}$.
\end{lemma}

\begin{proof}
Let $\rho=2C_t/\eta\ge2$, $\Lambda=\frac12\log X$ and $Y_k=\exp(\Lambda\rho^{-k})$ for $0\le k\le R$. If
$x\in[Y_k,Y_k^2]$ and $x^\eta<e\le x^{C_t}$, then $\log e\in I_k:=(\eta\Lambda\rho^{-k},2C_t\Lambda\rho^{-k}]
=(\eta\Lambda\rho^{-k},\eta\Lambda\rho^{1-k}]$. The $R+1$ intervals $I_0,\dots,I_R$ are pairwise disjoint, so some
$I_k$ contains no $\log e$ with $e\in\mathrm{Exc}$; take $Y=Y_k$. Then $Y^2\le Y_0^2=X$ and
$\log Y\ge\Lambda\rho^{-R}=\log X/(2^{R+1}(C_t/\eta)^R)\ge\log X/J$.
\end{proof}

\begin{theorem}[$2$-part and bad layers]\Fmark\label{thm:window:bab}
Assume Poitou--Tate duality, local finiteness and duality over finite extensions of $\Ql$, \cref{lit:window:fw}
and $\sfname{LemmaK}(C_K)$. There are functions $R_0,R_\ell$ on abelian data such that the following holds for all
$C_t\ge1$. Let $M$ be an abelian datum, $B'$ a finite set of odd primes with abelian data $L(\ell)$ ($\ell\in B'$)
whose fields contain $\mu_\ell$, $R=R_0(M)+\sum_{\ell\in B'}R_\ell(L(\ell))$, $\eps\in(0,\frac1{52})$,
$J=2^{R+2}(2C'_KC_t/\eps)^R$ with $C'_K=\max(C_K,1)$, and $A_1>0$. For all large $X$, all $m,t:\R\to\N$ with $2\le m(x)$,
$2^{m(x)}\le x^{A_1}$, $2\le t(x)\le x^{C_t}$, and all primes $q\equiv3\bmod8$, $d\equiv7\bmod16$ prime to $m_M$
with $\log(qd)\le X^{\eps/(4J)}$, there is $Y\in[X^{1/J},X^{1/2}]$ such that for all $x\in[Y,Y^2]$:
\begin{enumerate}[label=(\alph*)]
\item every $E=E_{M,m(x),q,d}$ satisfies $\sfname{VVat}(2,E,t(x),\exp(x^{4\eps/5}))$;
\item for $\ell\in B'$, every layer $L_1$ of degree $\ell^{j_\ell}$, $j_\ell=\lfloor\log_\ell R_\ell(L(\ell))\rfloor+3$,
  over a field of type $L(\ell)$ satisfies $\sfname{VVat}(\ell,L_1,t(x),\exp(x^{4\eps/5}))$.
\end{enumerate}
\formal{UpperTheory.V.Win.thmBab_of_lemmaK}
\end{theorem}

\begin{proofoutline}
By \cref{lit:window:fw,lem:window:tower}: (1) every abelian $L_0\ni\mu_\ell$, $\ell$ odd, satisfies
$\sfname{ABound}(\ell,R,L_0)$ for some $R$; (2) for every $M$ there is $R$ such that for all $(q,d)$ there is
$L_0\subset\Q(\zeta_{4m_Mqd})$ containing $\mu_4$, of degree $\le8\deg M$, with $\sfname{ABound}(2,R,L_0)$, and
$\windowsh(E,2,j)\le1+\windowsh(L_1,2,j)$ for every $E$ in (a) and a layer $L_1$ of the cyclotomic $\Z_2$-tower of
$L_0$ of index $\le m(x)+2$. Let $\mathrm{Exc}$ be the set of the at most $R$ exceptional twists $e\le X^{C_t}$ of the
relevant towers. Non-exceptional twists are bounded by the first estimate of $\sfname{ABound}$, with $k\ll\log x$.
If $x$ satisfies the conclusion of \cref{lem:window:window} with $\eta=\eps/(2C'_K)$, every exceptional $e\le t(x)$
is $\le x^{\eps/(2C'_K)}$, and the second estimate together with $\sfname{LemmaK}$ for $L_2$, whose discriminant is
$\le\exp(c_0x^{\eps/4})$, bounds these twists by $\exp(x^{4\eps/5})$ as well.

%P2PTR thm:window:bab
A complete proof is given in \cref{sec:pf:vwin}.
\end{proofoutline}

\begin{remark}[Iwasawa towers and Behrend's construction]\label{rem:window:iwasawa}
The announcement \cite{Announcement2026} motivates the Iwasawa tower by analogy with a pigeonhole argument over the
constant field extensions $\F_{q^n}/\F_q$ of a function field. In the formal proof the condition
$a\equiv\pm1\bmod2^{m+2}$ places $E$, up to index two, in the cyclotomic $\Z_2$-tower over a base of bounded degree
(step (2) above); the control along towers is \cref{lem:window:tower}, and the pigeonhole argument is
\cref{lem:window:window}. Behrend's construction \cite{Behrend1946}, mentioned in \cite{Announcement2026} for the
precise choice of the tower, appears in the construction of the hosts (\cref{sec:hosts}): there one needs, for
integers $n$ and $L\ge2$, sets $T\subset[0,n/(4L))$ of size at least $n\exp(-C\sqrt{\log n\log L})/L$ without relations
$j(s_2-s_1)=i(s_3-s_1)$ ($s_1,s_2,s_3\in T$ distinct, $0<|i|,|j|\le L$ distinct), and these are obtained from
Behrend's spheres. $\sfname{VWin}$ itself holds for all fields $E_{\psi,m,q,d}$.
\end{remark}

% ---------------------------------------------------------------------------
\subsection{The plus part}\label{sec:window:plus}

\begin{definition}[Hosts and character components]\label{def:window:host}
A \emph{host} $h=(\chi,m,q_1,d)$ consists of a primitive character $\chi$, $m\ge2$ and primes $q_1\equiv3\bmod8$,
$d\equiv7\bmod16$ prime to $N_\chi$. Let $N_h=\lcm(2^{m+2},N_\chi)q_1d$ and $\windowH_h\subset(\Z/N_h)^\times$ the
subgroup of \cref{def:window:fields}, with fixed field $E_h=E_{\chi,m,q_1,d}$; then $\cond(E_h)=N_h$. For a
coefficient ring $\OO$ (the ring of integers of a finite extension of $\Ql$ inside $\Qlbar$, with ramification index
$e_\OO$) and $j\ge0$, let $\windowX_h(\OO,j)$ be the set of characters $\theta$ of $(\Z/N_h)^\times/\windowH_h$ with
values in $\OO^\times$ and $\theta(-1)=(-1)^j$, and put
\[
 V_\ell(\theta,j)=\len_\OO\Sha_{\ell N_h}\bigl(\Q,\OO\otimes\windowQZ(\theta^{-1}\windowcyc^{-j})\bigr),\qquad
 \widetilde V_\ell(\theta,j)=V_\ell(\theta,j)/e_\OO,
\]
where $\Sha_n(\Q,M)\subset H^1(G_\Q,M)$ consists of the classes vanishing on all inertia groups and on the
decomposition groups at the primes dividing $n$.
\end{definition}

By Shapiro's lemma and inflation--restriction, for odd $\ell$, $j\ne0$ and $\OO$ with enough roots of unity,
$e_\OO\,\mathrm{plusLen}(D_\ell,j)\le\sum_{\theta\in\windowX_h(\OO,j)}V_\ell(\theta,j)$ if
$\ell\nmid[(\Z/N_h)^\times:\windowH_h]$, and in general up to an explicit loss term polynomial in $N_hj$.

\begin{definition}[Bad primes and cases]\label{def:window:cases}
For primitive $\chi$ let $p_K(\chi)$ be the least prime $p\equiv3\bmod4$, $p\ge7$, $p\nmid N_\chi$, let $h^*(\chi)$ be
the class number of $K^*_\chi=\Q(\sqrt{-p_K(\chi)})$, and let $p_1(\chi)$ be the least prime $\equiv1$ modulo the odd
part of $o_\chi$ and prime to $2N_\chi p_K(\chi)$. The \emph{bad set} is
$B_\chi=\{3\}\cup\{\ell\text{ odd}:\ell\mid o_\chi N_\chi p_K(\chi)h^*(\chi)p_1(\chi)\}$. A host is \emph{big} for $t$ if
$2t+1<q_1,d$ and $N_\chi,o_\chi,p_K(\chi),h^*(\chi),p_1(\chi)<q_1,d$. For odd $\ell$, $0<j<t$ and
$\theta\in\windowX_h(\OO,j)$, let $\theta^\flat$ be the prime-to-$\ell$-order part of $\theta$ with its
$\ell$-component removed. The cases are listed in \cref{tab:window:cases}; if $h$ is big for $t$, every triple
$(\ell,j,\theta)$ falls into one of them.
\end{definition}

\begin{table}[ht]
\centering\small
\begin{tabular}{@{}llll@{}}
\toprule
case & condition on $(\ell,j,\theta)$ & bound & method \\
\midrule
I & $\ell\notin B_\chi$, $(\ell-1)\nmid2j$ & $\sum_\ell$: $(N_ht)^C$ & $\Uni(3)$; \cref{thm:kato}\\
Z & $\ell\notin B_\chi$, $(\ell-1)\mid2j$, $\theta=1$ & $V_\ell(1,j)=0$ & \cref{prop:window:propZ}\\
II & $j\ge2$, $\ord\theta^\flat$ even, $\ell\in B_\chi$ or $(\ell-1)\mid2j$ & $(\ell N_hj)^C$ & $\GSp_4$\\
III & $\ell\notin B_\chi$, $(\ell-1)\mid2j$, $\ord\theta$ odd, $\theta\ne1$ & $(\ell N_hj)^C$ & $\GSp_4$\\
IV & $\ell\in B_\chi$, $j\ge2$, $\ord\theta^\flat$ odd & layers${}+(N_ht)^C$ & \cref{thm:window:bab}(b)\\
V & $\ell\in B_\chi$, $j=1$ & $(\ell N_hj)^C$ & $\Uni(3)$ over $\Q(i)$\\
\bottomrule
\end{tabular}
\caption{The case division for the plus part; the bounds are for $\widetilde V_\ell(\theta,j)\log\ell$.}\label{tab:window:cases}
\end{table}

\begin{definition}[Counts]\label{def:window:counts}
For a case $P$, $\sfname{Count}P$ says: for every primitive $\chi$ there is $C$ such that for every host $h$ over
$\chi$ big for $t$, every odd $\ell$, every $\OO$, every $0<j<t$ and every $\theta\in\windowX_h(\OO,j)$ in case $P$,
$V_\ell(\theta,j)<\infty$ and $\widetilde V_\ell(\theta,j)\log\ell\le(\ell N_hj)^C$. This defines $\sfname{CountII}$,
$\sfname{CountIII}$, $\sfname{CountV}$ (case V under an auxiliary $2$-adic condition) and $\sfname{CountZ}$.
$\sfname{CountI}$ is summed over primes: for every $\chi$ there are $C$ and primes $r_1,r_2$ such that for hosts with
$r_1,r_2<q_1,d$, big for $t$, all $0<j<t$, all coefficient rings and every finite set $L$ of odd primes in case I,
$\sum_{\ell\in L}\log\ell\sum_{\theta}\widetilde V_\ell(\theta,j)\le(N_ht)^C$.
\end{definition}

\begin{proposition}\Fmark\label{prop:window:propZ}
For $\ell\ge5$ and even $j\ge2$ with $(\ell-1)\mid2j$, $\Sha_\ell(\Q,\OO\otimes\windowQZ(-j))=0$. Hence $\sfname{CountZ}$
holds.
\end{proposition}
%P2PTR prop:window:propZ
\begin{proof}See \cref{sec:pf:vwin}.
\end{proof}


If $(\ell-1)\mid j$, the residual module is trivial, and a control argument reduces the statement to the absence of
cyclic extensions of $\Q$ of degree $\ell$ unramified outside $\ell$ and split above $\ell$ (Kronecker--Weber); if
$j\equiv\frac{\ell-1}2\bmod(\ell-1)$, the residual character is the quadratic character of the quadratic subfield of
$\Q(\mu_\ell)$, and an explicit computation with that field applies. Local conditions at $N_h$ only shrink the group;
so $V_\ell(1,j)=0$ in case Z, where $\ell\ge5$ as $3\in B_\chi$, and $j$ is even as $\theta=1$.

\begin{theorem}[Plus part]\Fmark\label{thm:plus-part}
Assume Poitou--Tate duality, local finiteness and duality over finite extensions of $\Ql$, $\sfname{LemmaK}(C_K)$ and
$\sfname{CountI}$, $\sfname{CountII}$, $\sfname{CountIII}$, $\sfname{CountV}$, $\sfname{CountZ}$, and let $C_t>0$. For
every primitive $\psi$ there are $C$ and primes $r_1,r_2$ with the following property. Let $\eps\in(0,\frac1{52})$,
$A_1>0$, and let $J_\psi$ be the constant $J$ of \cref{thm:window:bab} for the data $M_\psi$, $B'=B_\psi$ and
$L(\ell)=\Q(\zeta_{\ell N_\psi})$. For all large $X$, all $m,t$ as in \cref{thm:window:bab}, all primes
$q\equiv3\bmod8$ and $d\equiv7\bmod16$ exceeding $2X^{C_t}+1$, $N_\psi$, $o_\psi$, $p_K(\psi)$, $h^*(\psi)$,
$p_1(\psi)$, $r_1$ and $r_2$, with $\log(qd)\le X^{\eps/(4J_\psi)}$, and every $Y$ satisfying the conclusion of
\cref{thm:window:bab} for these data, every field $E=E_{\psi,m(x),q,d}$ with $x\in[Y,Y^2]$ satisfies
\[
 \sfname{VVplusLe}\Bigl(E,t(x),\bigl(\cond(E)\,t(x)\bigr)^C+C\exp(x^{4\eps/5})\Bigr).
\]
\formal{UpperTheory.V.HV.thmHV_of_counts}
\end{theorem}

\begin{proofoutline}
\emph{Good primes.} For odd $\ell\notin B_\psi$, every $(\ell,j,\theta)$ is in case I, Z, II or III. Case I is
summed over $\ell$ by $\sfname{CountI}$. In the other cases $(\ell-1)\mid2j$, so $\ell\le2t$, and the bounds
$(\ell N_hj)^C\le(N_ht)^{3C}$ can be summed over $\ell\le2t$, $j<t$ and the at most $N_h$ characters; for these $\ell$
the reduction to characters has no loss term. \emph{Bad primes.} For each of the finitely many $\ell\in B_\psi$,
$\log\ell\sum_{j<t}\mathrm{plusLen}(D_\ell,j)\le(N_ht)^C+CB$ as soon as the layers $L_1$ of $\Q(\zeta_{\ell N_\psi})$
satisfy $\sfname{VVat}(\ell,L_1,t,B)$: case IV by realizing $\theta$ on such a layer; case V by $\sfname{CountV}$ or,
under the complementary $2$-adic condition, by Poitou--Tate duality, $\windowsh\le\mathrm{VLen}$ and $\sfname{LemmaK}$
on a field of bounded degree containing $\mu_\ell$. \emph{Assembly.} Take $B=\exp(x^{4\eps/5})$, as provided by
\cref{thm:window:bab}(b), and note that $N_h=\cond(E)$.

%P2PTR thm:plus-part
A complete proof is given in \cref{sec:pf:vwin}.
\end{proofoutline}

\begin{remark}[Plus parts and Vandiver's conjecture]\label{rem:window:vandiver}
In \cite{Announcement2026} this part of the argument is described as ``a significant weakening of Vandiver''
for those classes in $H^1(\OO_E[1/\ell],\windowQZ(-e))$ which cannot be killed by twisting. The formal
statements support the following reading. For $\ell\nmid[E:\Q]$ the trivial-character component of $\Sha_D(j)$ is
$\Sha^1(G_{\Q,S_\ell},\windowQZ(-j))$. For even $j$ it lies in the plus part and is dual to a subgroup of
$H^2(\Z[1/\ell],\Zl(1+j))$, which is isomorphic to $K_{2j}(\Z)\otimes\Zl$ \cite{Soule1979,Voevodsky2011}; and
Vandiver's conjecture for $\ell$ is equivalent to the vanishing of $H^2(\Z[1/\ell],\Zl(1+j))$ for all even $j\ge2$
(see \cite{Kurihara1992}). On plus components the values $L(\theta,-j)$ vanish and give no control. Except in case Z,
$\sfname{VWin}$ does not assert vanishing but explicit polynomial bounds (\cref{tab:window:cases}), uniformly over
all plus components of the fields $E$.
\end{remark}

% ---------------------------------------------------------------------------
\subsection{Eisenstein congruence numbers for \texorpdfstring{$\Uni(3)$}{U(3)} and \texorpdfstring{$\GSp_4$}{GSp(4)}}\label{sec:window:eis}

The counts are proved by one pattern, a quantitative form of Ribet's method \cite{Ribet1976} in the higher-rank
setting of \cite{BellaicheChenevier2009}:
\begin{enumerate}[label=(\roman*)]
\item $\Sha_{\ell N}(\Q,\OO\otimes\windowQZ(\theta^{-1}\windowcyc^{-j}))$ is compared with tangent spaces of a framed
  deformation problem for a residually reducible representation, of rank $3$ and unitary type over an auxiliary
  imaginary quadratic field $K$, or of rank $4$ and symplectic type, with Eisenstein semisimplification built from
  $\Psi=\windowcyc^j\theta$ (in rank $3$, $\Psi\oplus1\oplus\Psi^{-1}$);
\item commutative algebra bounds the length by twice the sum, over the characteristic-zero points $r_i$ of the
  deformation ring, of $v_\ell(\tr r_i(\sigma)-\tr r_{\mathrm{Eis}}(\sigma))$ for a chosen Frobenius element $\sigma$;
\item automorphy and rigidity statements show that the $r_i$ are automorphic, of bounded level and weight;
\item an \emph{Eisenstein congruence number}, a nonzero integer whose logarithm is bounded polynomially in the level
  and the weight, is divisible by the product of these differences;
\item summing over $\ell$ with weights $\log\ell$ is bounded by its logarithm; two auxiliary primes $p\ne p'$, at
  which the Frobenius elements are taken, cover the case $\ell=p$.
\end{enumerate}

\emph{Case I} uses $K=K^*_\chi$. The statement $\sfname{PropGK}$ asserts: for every imaginary quadratic $K$ there are
$c_5,A$ such that for all $j\ge1$, $N$, subgroups $H\subset(\Z/N)^\times$ of index $i_H$ and primes $p\nmid30Nd_K$
there is an integer $\mathcal N\ne0$ with $\log|\mathcal N|\le c_5\varphi(i_H)(N(j+1))^A\log p$ such that for all
primes $\ell\ne p$ and all $\OO$: if some character with kernel $H$ satisfies a list $\sfname{GD}(K,\ell,j,\cdot)$ of
genericity conditions (among them $(\ell-1)\nmid2j$ and $\ell\nmid d_Kh_K$), then
$\sum_{\ker\theta=H}\widetilde V_\ell(\theta,j)\le2v_\ell(\mathcal N)$, all terms being finite. It follows from a
\emph{number} statement ($v_\ell(\mathcal N)$ dominates the sum of $v_\ell(\tr\rho(\sigma_w^{-1})-\mathrm{eis}_3)$ over every
finite set of pairwise non-isomorphic $3$-dimensional representations $\rho$ of $G_K$ of bounded level which are of
unitary type or endoscopic, i.e.\ contain the representation $W$ of a newform of weight $2j+1$; here $\sigma_w$ is
a Frobenius element at $w\mid p$) and a \emph{one-character} statement (under $\sfname{GD}$,
$\widetilde V_\ell(\theta,j)\le2\sum_{\rho\in X}v_\ell(\tr\rho(\sigma)-\Psi(\sigma)-1-\Psi(\sigma)^{-1})$ for such a finite set
$X$).

\begin{theorem}[Case I]\Fmark\label{thm:window:countI}
$\sfname{HypothesesH'}$, $\sfname{LemmaK}(C_K)$, $\sfname{KatoTwoW}$ (\cref{def:window:kato}) and an automorphy
statement of $R=T$ type for the $3$-dimensional deformation problems with unramified coefficients imply
$\sfname{CountI}$.
\formal{UpperTheory.V.Z.J.count_i_of_printed_bound_k_kato_automorphy}
\end{theorem}
%P2PTR thm:window:countI
\begin{proof}See \cref{sec:pf:eis}.
\end{proof}


The one-character statement uses $\sfname{LemmaK}$, Poitou--Tate duality, the Greenberg--Wiles formula, a
presentation lemma for framed deformation rings, rigidity statements derived from an adjoint lifting theorem, and
the automorphy statement. The vanishing theorem $\sfname{KatoTwoW}$ enters at the endoscopic points, which are
block embeddings of a $2$-dimensional $W$: their adjoint representation contains $W$ and
$W\otimes\vartheta\cong W^\vee$, whose Bloch--Kato Selmer groups vanish by \cref{thm:kato}.

\emph{Cases II and III} use $4$-dimensional representations of $G_\Q$ of Hodge--Tate type $(j+1,j,1,0)$ with
Eisenstein traces of $\GSp_4$-type. The automorphy of all irreducible $4$-dimensional points of the relevant
deformation problems is deduced from $\sfname{HypothesesH'}$ by potential modularity in towers of totally real and CM
fields, using ordinary Hilbert modular forms congruent to Eisenstein series (with a weight-changing congruence
theorem proved by a trace formula argument), automorphic induction (\cref{rem:window:bc}(iii)) and an unramified
adjoint lifting theorem.
\emph{Case V} ($j=1$, $\ell\in B_\chi$) applies the unitary pattern over $\Q(i)$, with an automorphy statement deduced
from the density of automorphic points for unitary groups in three variables, which uses $\sfname{Row0b}$.

% ---------------------------------------------------------------------------
\subsection{Vanishing of Bloch--Kato Selmer groups}\label{sec:window:kato}

Let $\ell$ be odd. For continuous $W:G_\Q\to\GL_2(\Qlbar)$, $\windowHf(\Q,W)=0$ means that every continuous cocycle
which is unramified at all primes $\ne\ell$ and crystalline at $\ell$ is a coboundary \cite{BlochKato1990}.

\begin{definition}[$\sfname{KatoTwoW}$]\label{def:window:kato}
Let $j\ge1$ with $(\ell-1)\nmid2j$, $K$ an imaginary quadratic field with quadratic character $\vartheta$, $\psi$ a
character of finite order prime to $\ell$, and $W:G_\Q\to\GL_2(\Qlbar)$. The configuration $\sfname{Used}$ requires,
among other conditions: $\det W=\vartheta$; $W$ has finite level away from $\ell$; $W$ is ordinary at $\ell$, i.e.\ upper
triangular on a decomposition group at $\ell$ with diagonal characters $\windowcyc^{j}\chi'$ and $\windowcyc^{-j}\chi'$
on inertia, where $\chi'$ has order dividing $\ell-1$; $W$ is irreducible on $G_{K(\mu_{\ell^\infty})}$; and
$\tr W\equiv\vartheta\Psi+\Psi^{-1}$ residually, where $\Psi=\windowcyc^j\psi$. $\sfname{KatoTwoW}$ asserts: for every
such $W$, every $\iota:\Qlbar\cong\C$ and
every newform $f$ of weight $2j+1$ attached via $\iota$ to $W\otimes\windowcyc^{-j}$, if $W$ satisfies the window
condition $\sfname{WitAQ}(\iota,W)$ (for imaginary quadratic $K'$ with $W|_{G_{K'}}$ irreducible, the Satake
parameters of the cuspidal representations of $\GL_2(\A_{K'})$ attached to $W|_{G_{K'}}$ lie in windows of width
one), then $\windowHf(\Q,W)=0$ and $\windowHf(\Q,W\otimes\vartheta)=0$.
\end{definition}

\begin{theorem}[Kato-type vanishing]\Fmark\label{thm:kato}
$\sfname{HypothesesH'}$ implies $\sfname{KatoTwoW}$.
\formal{UpperTheory.V.Z.Kato.WeilSphericalRecords.katoTwoW_of_printsH}
\end{theorem}
%P2PTR thm:kato
\begin{proof}See \cref{sec:pf:kato}.
\end{proof}


As $f$ has odd weight, all critical values of its twists are non-zero (no integer lies at the centre;
Jacquet--Shalika \cite{JacquetShalika1976} at the edge; functional equation), and the theorem is in substance a
special case of \cite[Thm.~14.2]{Kato2004} for ordinary forms and the two near-central twists $W$ and
$W\otimes\vartheta\cong W^\vee$. The formal proof re-derives it along Kato's method, without assuming Kato's theorem.
Only part of $\sfname{Used}$ is needed (not, e.g., the residual condition on $\tr W$), and $\sfname{WitAQ}$ serves
only to show that the automorphic representations attached to twists of $W$ are unramified at $\ell$.

\begin{proofoutline}
\emph{Reduction.} $W$ is \emph{admitted}: continuous, $\det W$ quadratic and odd, irreducible on
$H_0=\ker\windowcyc\cap\ker\det W$, of finite level away from $\ell$, modular and ordinary. Local--global
compatibility at $\ell$ for $\GL_2$ over $\Q$, derived from $\sfname{Row0b}$ by base change to an imaginary
quadratic field in which $\ell$ splits and from the local parameters of \cref{sec:galois}, shows together with
$\sfname{WitAQ}$ that the automorphic representations attached to twists of $W$ are unramified at $\ell$. This
yields an $\ell$-ordinary newform $g$ of level prime to $\ell$ attached to a twist of $W$, whose unit root is not a
root of unity. As $W\otimes\vartheta=W\otimes\det W$ is again admitted, it suffices to show $\windowHf(\Q,W)=0$.

\emph{The Euler system.} After conjugation $W$ is integral, unramified outside a finite set $S\ni\ell$ and upper
triangular at a place above $\ell$. One constructs a bottom class $y\in Z^1(G_\Q,W)$, unramified away from $\ell$,
whose image in the quotient line at $\ell$ is not a coboundary, and for all $n$, all primes $w\notin S'$ with
$\ell^n\mid w-1$ and all surjective $\kappa:G_\Q\to\Z/\ell^n$ of conductor $w$, a class over the field cut out by
$\kappa$ with corestriction $P_w\,y$, where $P_w\equiv u\det(W(\Frob_w)-1)\bmod\ell^n$ for a unit $u$. The construction
follows Kato: zeta elements from Siegel units \cite{KubertLang1981}, with the Kummer theory of modular units in the
field of modular functions of all levels \cite{Shimura1971} replacing \'etale cohomology of modular curves;
projection to the $g$-part via an Eichler--Shimura decomposition; norm relations; an explicit reciprocity law
computing the Coates--Wiles derivatives along $\Q_\ell(\mu_{\ell^\infty})$ by Gauss sums of an Eisenstein measure
\cite{CoatesWiles1977,Coleman1979}; and the non-triviality of the bottom class at $\ell$, from Coleman power series and
the non-vanishing of $L(g\otimes\phi,s)$ at the edge of the critical strip.

\emph{From the Euler system to vanishing.} Let $W^*=W\otimes\windowcyc(\det W)^{-1}\cong W^\vee(1)$. The strict Selmer
group of $W^*$ (classes unramified away from $\ell$ and locally trivial at $\ell$) vanishes by
\cref{lem:window:heights}; by local Tate duality for characters so does its Greenberg Selmer group; by the
Greenberg--Wiles formula so does the Greenberg Selmer group of $W$; and every Bloch--Kato class of $W$ satisfies the
Greenberg condition, since the quotient character at $\ell$ has Hodge--Tate weight $-j<0$.
\end{proofoutline}

The Euler-system argument uses a single auxiliary prime at a time, whereas the classical arguments
\cite{Kolyvagin1990,Rubin2000} pair derivative classes with dual classes at many auxiliary primes simultaneously. Let
$\omega(x,y)=x_0y_1-x_1y_0$, so that
$\omega(W(s)x,W^*(s)y)=\windowcyc(s)\omega(x,y)$. For integral cocycles $y$ of $W$ and $z$ of $W^*$ with
$y\cup z\equiv du\pmod{\ell^n}$ (possible when $z$ is strict, as $\Sha^2(G_S,\mu_{\ell^n})=0$) put
$\mathrm{ht}(s)=\det(W(s)-1)u(s)-\omega\bigl(\mathrm{adj}(W(s)-1)y(s),z(s)\bigr)$.

\begin{lemma}[Heights on commutators]\Fmark\label{lem:window:heights}
In this situation, assume the tame relations of the Euler system for $y$. Then
$\det(W(\sigma)-1)\,\mathrm{ht}(\sigma)\equiv0\pmod{\ell^n}$ for every $\sigma\in[H_0,H_0]$. Consequently the strict
Selmer group of $W^*$ vanishes.
\end{lemma}

\begin{proofoutline}
By Chebotarev's theorem choose a prime $w\notin S\cup S'$ with $\ell^n\mid w-1$, a Frobenius element $\tau$ at $w$
close to $\sigma$, and a character $\kappa$ of conductor $w$ trivial at the places of $S$. The tame relation gives a
cochain $d$ (Kolyvagin's derivative) with $d(st)\equiv d(s)+W(s)d(t)+\kappa(s)P_wW(s)y(t)$. The $2$-cocycle
$\xi=d\cup z-\kappa P_w\windowcyc u$ is a local coboundary at all places of $S$, hence, by the reciprocity law, also
at $w$, and a local computation at $w$ gives $P_w\,\mathrm{ht}(\tau)\equiv0$, i.e.\
$\det(W(\tau)-1)\,\mathrm{ht}(\tau)\equiv0$; by continuity the same holds at $\sigma$. For the consequence, let $z$ be a
non-trivial strict class. Then $y$ and $z$ do not vanish on $N_0=\ker\windowcyc\cap\ker W$ (otherwise $y$ would
satisfy the Greenberg condition, resp.\ $z$ would be a coboundary), and there are $g,h\in N_0$ with
$\delta'=\omega(y(h),z(g))-\omega(y(g),z(h))\ne0$ and $\phi\in[H_0,H_0]$ with $\delta=\det(W(\phi)-1)\ne0$. A direct
computation gives $\mathrm{ht}(\phi[g,h])-\mathrm{ht}(\phi)\equiv\delta\delta'$, so the first part at $\phi$ and
$\phi[g,h]$ yields $\delta^2\delta'\equiv0\pmod{\ell^n}$ for all $n$, a contradiction.

%P2PTR lem:window:heights
A complete proof is given in \cref{sec:pf:kato}.
\end{proofoutline}

In the formal development the crystalline condition at $\ell$ refers to a period ring characterized by a list of
axioms; it is used only through Bloch--Kato's local dimension formula.
